% Generated draft: requires mathematical and source review.
\documentclass[11pt]{article}
\usepackage[margin=25mm]{geometry}
\usepackage{fontspec}
\setmainfont{Latin Modern Roman}
\usepackage{amsmath,amssymb,mathtools,unicode-math}
\setmathfont{Latin Modern Math}
\usepackage{enumitem,needspace,xurl,hyperref}
\setlength{\parindent}{0pt}
\setlength{\parskip}{5pt}
\setlength{\emergencystretch}{3em}
\newcommand{\N}{\mathbb{N}}
\newcommand{\Z}{\mathbb{Z}}
\newcommand{\Q}{\mathbb{Q}}
\newcommand{\R}{\mathbb{R}}
\newcommand{\C}{\mathbb{C}}
\newcommand{\F}{\mathbb{F}}
\newcommand{\A}{\mathbb{A}}
\newcommand{\PP}{\mathbb{P}}
\newcommand{\E}{\mathbb{E}}
\newcommand{\eps}{\varepsilon}
\newcommand{\dd}{\,\mathrm{d}}
\begin{document}
\section*{UnsolvedMath ID: 3100038}
\textbf{Source problem code:} catalogue entry
\subsection{Definitions and mathematical statement}
The source asks three enumeration questions: determine the Whitney numbers of the lattice of contractions of the $n$-cube; determine the counts when cube-symmetric contractions are identified; and count contractions up to isomorphism. The precise cube/contraction model is not defined.

Under the graphic interpretation suggested by the summary, let $Q_n$ have vertices $\{0,1\}^n$ with adjacency at Hamming distance one. Its graphic-matroid flat lattice can be described by partitions $\pi$ of the vertices whose blocks induce connected subgraphs, ordered by refinement. An associated contraction collapses each block. The rank in this lattice is $r(\pi)=2^n-|\pi|$.

For a finite ranked lattice $L$ with bottom $0$, the Whitney numbers of the second kind are
\[
W_j(L)=\#\{x\in L:r(x)=j\},
\]
and signed Whitney numbers of the first kind are
\[
w_j(L)=\sum_{r(x)=j}\mu_L(0,x),
\]
where $\mu_L$ is the M\"obius function. The source does not specify which kind is requested.

For this graph model, cube symmetries act by coordinate permutations and coordinate flips. A natural symmetry-counting target is the number of their orbits in each rank. Such an orbit count is not automatically the first-kind Whitney number of a quotient lattice; no quotient-lattice order is supplied.

Finally, isomorphism counting could concern the quotient graphs obtained by collapsing blocks, rather than orbits under symmetries of the original cube. These equivalences are distinct. Retaining parallel edges or replacing the quotient by a simple graph changes the count, as does using matroid rather than graph isomorphism. Moreover, the source's $n$-cube could refer to a different geometric/matroidal object, so the graphic connected-partition model is an explicit candidate, not a silently reconstructed definition. The questions require that model and the requested Whitney/isomorphism conventions before numerical values can be asserted.
\subsection{Short English statement}
Enumerate cube contractions by rank, by cube-symmetry orbits, and by isomorphism type. The contraction model, Whitney-number kind and treatment of quotient graphs are not specified in the source.
\subsection{Sources}
\begin{itemize}
\item Local UnsolvedMath registry, permanent ID 3100038, exact source code catalogue entry; no proposer or year supplied. Registry status open, classification OPEN-TRIAGE.
\item Joshua Cooper, \emph{Combinatorial Problems I Like}, October 2020 snapshot, Matroids, Lattices, and Posets, source-order bullet 38, \url{https://people.math.sc.edu/cooper/combprob.html}. Supplied HTML extraction accessed July 29, 2026. The source's additional ``See this'' link target is absent.
\item The research summary suggests the graphic matroid and mentions computational/small cases and vague quotient or shard literature without exact references. It calls the item partly addressed/literature-survey, but does not define the lattice or establish exact arbitrary-$n$ counts.
\item Provenance flags status and rights NEEDS\_REVIEW: the 2020 snapshot labels the item open, later resolution was not checked item by item, and public access did not establish redistribution permission.
\end{itemize}
\end{document}
